% ch13 -- Alles ist Information --- aber welche?
% Quelldokumente: A180, 243, 247

\epigraph{Wenn man alle Einheiten auf 1 setzt, bleiben Verhältnisse.
Reine Zahlen. Strukturen. Information.}{Matrix-Buch, Kap.\,16}

\section{Die vierte Lesart}

Die T0-Theorie kennt drei gleichberechtigte Lesarten: Alles ist Energie,
alles ist Masse, alles ist Zeit. Es gibt eine vierte:

\textbf{Alles ist Information.}

Was bleibt invariant, wenn man alle Einheiten auf 1 setzt
($c = \hbar = \alpha = G = 1$)? Es bleiben \textbf{Verhältnisse}.
Reine Zahlen. Strukturen. Das Massenverhältnis $m_e/m_\mu \approx 1/206{,}8$
ist keine Energie --- es ist eine Relation, ein Stück reiner Information
über die Geometrie des Torus. \status{B}

\section{Information ohne Träger?}

Hier liegt die entscheidende Spannung zu A271--A273:

In der Landauer-Analyse ist Information \emph{immer} an einen Träger
gebunden. Ein Bit ohne Träger ist eine abstrakte Kategorie, keine
physikalische Größe --- und hat keine Energie. \status{B}

In der kosmologischen Lesart sind die Verhältnisse $\xi, m_e/m_\mu, G/\hbar c$
\emph{selbsttragende} Informationen --- sie sind nicht in einem äußeren
Medium kodiert. Das Universum ist sein eigener Träger.

Das ist kein Widerspruch, sondern eine Grenzfall-Situation: Wenn das Medium
\emph{selbst} aus Information besteht, gibt es keinen unabhängigen Träger
mehr.

\begin{laierspur}
Ein Buch enthält Information auf Papier --- der Träger ist das Papier.
Aber was ist, wenn das Papier selbst aus Buchstaben besteht?
Dann ist der Träger und der Inhalt dasselbe. Das ist die Situation im
Universum, wenn man es als geometrische Struktur versteht: Die Geometrie
ist das Trägermedium und die Information gleichzeitig.
\end{laierspur}

\section[\texorpdfstring{$\xi$}{xi} aus der Galois-Algebra]{$\xi$ aus der Galois-Algebra}

In der FFGFT ist $\xi = 4/30000$ keine offene Frage. Sie folgt zwingend
aus den Galois-Gruppenordnungen des Körpers $\mathrm{GF}(3^n)$ und den
Wicklungszahlen des $T^4$-Torus (Dok.\,338):
\[
  \xi
  = \frac{(r_\mu/r_e)^2}{|\mathrm{GF}(9)^*|^2 \cdot 5^2 \cdot |\mathrm{GF}(27)|}
  = \frac{(12/5)^2}{8^2 \cdot 25 \cdot 27}
  = \frac{4}{30000}
\]
Die drei Galois-Faktoren haben konkrete algebraische Bedeutung: \status{K Dok.\,338}

$8 = |\mathrm{GF}(9)^*| = 3^2 - 1$ (multiplikative Gruppe der 2.\,Generation),
$27 = |\mathrm{GF}(27)| = 3^3$ (Körper der 3.\,Generation),
$5 = $ kleinster Primteiler von $|\mathrm{GF}(81)^*| = 80$.

Zusätzlich ist $\xi$ empirisch verankert (Matrix-Buch, Kap.\,2):
$\xi = \alpha / (E_0/\mathrm{MeV})^2$.

Die Antwort auf das „Warum" von $\xi$ ist damit algebraisch, nicht
spekulativ. In Matzkes Rahmen bleibt die analoge Frage --- Warum
$e_i^2 = +1$? --- tatsächlich offen.

\section{Wo Hyperbit und T0 konvergieren}

Matzkes Hyperbit: \enquote{Es gibt unterscheidbare Dinge} $\to$ Geometrie.

FFGFT: \enquote{Es gibt Windungen auf $T^4$} $\to$ Geometrie, Physik, Massen.

Beide starten mit reiner Unterscheidbarkeit --- reiner Information ---
und enden bei Physik. Das ist der tiefste gemeinsame Punkt beider Ansätze.
\status{B}
